% !PRINTED 697
\weberpage{697}
%<<CONT
\weq{5}{
\begin{aligned}
(O, B) &= (r_{s+1} - 2) + (r_{s+2} - 2) + \cdots + (r_n - 2), \\
p &= (r_{s+1} - 1) + (r_{s+2} - 1) + \cdots + (r_n - 1), \\
(O, B) &= p - n + s,
\end{aligned}
}
und nach (2):
\[
(O, A) - (O, B) = n - p + 1.
\]

Nun ist nach \S\,192, (14) $(O, B) = (A, W)$, und wir haben
den \gesperrt{Riemann-Rochschen} Satz für Klassen erster Gattung:

4. \gesperrt{Die Dimension einer eigentlichen Klasse $A$
erster Gattung ist}
\weq{6}{(O, A) = n - p + 1 + (A, W).}

Diese Formel schlie\ss t (3) in sich, weil für Polygone zweiter
Gattung $(A, W) = 0$ ist.

Da die Dimension einer eigentlichen Klasse mindestens
gleich 2 ist, so folgt für eine Klasse zweiter Gattung:
\[
n \geqq p + 1,
\]
und hierin ist der Beweis eines Satzes von \gesperrt{Riemann} enthalten:

5. \gesperrt{Eine Funktion, deren Ordnung kleiner als
$(p + 1)$ ist, ist von der ersten Gattung.}

Es folgt weiter daraus, da\ss\ die Klasse $W$ der Polygone
erster Gattung eine eigentliche ist. Denn angenommen, $W$ habe
den Teiler $\mathfrak{M}$, und es sei $W = \mathfrak{M} A$, so wäre $(O, W) = (O, A)$
$= p$. Nimmt man $\mathfrak{A}$ in $A$ relativ prim zu $\mathfrak{M}$, so gibt es in $W$ nur
ein durch $\mathfrak{A}$ teilbares Polygon, nämlich $\mathfrak{A}\mathfrak{M}$, und folglich ist
$(A, W) = 1$. Demnach gibt die Formel (6):
\[
n = 2p - 2,
\]
also gleich der Ordnung von $W$. Die Ordnung von $A$ ist also
ebenso gro\ss\ wie die von $W$, und folglich $\mathfrak{M} = \mathfrak{O}$.

Der Satz 4. gilt unverändert auch für uneigentliche Klassen.
Um das nachzuweisen, sei zunächst $A$ von der ersten Gattung
und vom Teiler $\mathfrak{M}$:
\weq{7}{
\begin{aligned}
A &= \mathfrak{M} A', \\
AB &= W, \\
A'B' &= W, \\
B' &= \mathfrak{M} B, \quad (\S\,192,\,8.) \\
(A', W) &= (O, B').
\end{aligned}
}

Die Ordnungen von $\mathfrak{M}, A, B$ seien $m, a, b$, also $a + b$ die
Ordnung von $W$, d. h.
\weq{8}{a + b = 2(p - 1).}
